\begin{kbvrule}
\noindent{\kbvsecfont ವ್ಯಾಖ್ಯೆ :}\hspace{0.5em}ಉತ್ತರಪದದ ಅರ್ಥವೇ ಪ್ರಧಾನವಾಗಿದ್ದು, ಪೂರ್ವ ಹಾಗೂ ಉತ್ತರಪದಗಳು ಲಿಂಗ ವಚನ ವಿಭಕ್ತಿಗಳಿಂದ ಸಮಾನವಾಗಿದ್ದು, ಅವೆರಡೂ ಪದಗಳು ವಿಶೇಷಣ ವಿಶೇಷ್ಯ ಸಂಬಂಧದಿಂದ ಕೂಡಿ ಆಗುವ ಸಮಾಸಕ್ಕೆ ಕರ್ಮಧಾರಯ ಸಮಾಸ ಎಂದು ಹೆಸರು.
\end{kbvrule}

\vspace{3pt plus 2pt}\noindent\hspace*{1.2em}\kbvlabel{ಉದಾ}\par\vspace{1pt}

\begin{kbvderiv}
\vspace{4pt plus 3pt minus 1pt}
\begin{center}\begin{adjustbox}{max width=\textwidth}
{\renewcommand{\arraystretch}{1.22}\setlength{\tabcolsep}{4pt}
\begin{tabular}{l c l}
ಹಿರಿದು + ಬಾಗಿಲು & {\color{kbvgrey}$=$} & \kbvresult{ಹೆಬ್ಬಾಗಿಲು}{\color{kbvgrey} {[}ವಿ.ಪೂ.ಕ{]}} \\
ಹಿರಿದು + ಜೇನು & {\color{kbvgrey}$=$} & \kbvresult{ಹೆಜ್ಜೇನು}{\color{kbvgrey} {[}ವಿ.ಪೂ.ಕ{]}} \\
ತಣ್ಣನೆ + ಗಾಳಿ & {\color{kbvgrey}$=$} & \kbvresult{ತಂಗಾಳಿ}{\color{kbvgrey} {[}ವಿ.ಪೂ.ಕ{]}} \\
ಸಿಹಿಯಾದ + ನೀರು & {\color{kbvgrey}$=$} & \kbvresult{ಸಿಹಿನೀರು}{\color{kbvgrey} {[}ವಿ.ಪೂ.ಕ{]}} \\
ಹಣ್ಣಾದ + ಎಲೆ & {\color{kbvgrey}$=$} & \kbvresult{ಹಣ್ಣೆಲೆ}{\color{kbvgrey} {[}ವಿ.ಪೂ.ಕ{]}} \\
ಎಳದು + ಕರು & {\color{kbvgrey}$=$} & \kbvresult{ಎಳಗರು}{\color{kbvgrey} {[}ವಿ.ಪೂ.ಕ{]}} \\
ಕಿರಿದು + ದಾರಿ & {\color{kbvgrey}$=$} & \kbvresult{ಕಿರುದಾರಿ}{\color{kbvgrey} {[}ವಿ.ಪೂ.ಕ{]}} \\
ಮೆಲ್ಲಿತು + ನುಡಿ & {\color{kbvgrey}$=$} & \kbvresult{ಮೆಲ್ನುಡಿ}{\color{kbvgrey} {[}ವಿ.ಪೂ.ಕ{]}} \\
ಕೆಚ್ಚನೆ + ಚೇಳು & {\color{kbvgrey}$=$} & \kbvresult{ಕೆಂಚೇಳು}{\color{kbvgrey} {[}ವಿ.ಪೂ.ಕ{]}} \\
ಇನಿದು + ಸರ & {\color{kbvgrey}$=$} & \kbvresult{ಇಂಚರ}{\color{kbvgrey} {[}ವಿ.ಪೂ.ಕ{]}} \\
ನೀಲವಾದ + ಆಕಾಶ & {\color{kbvgrey}$=$} & \kbvresult{ನೀಲಾಕಾಶ}{\color{kbvgrey} {[}ವಿ.ಪೂ.ಕ{]}} \\
ಶ್ವೇತವಾದ + ಅಂಬರ & {\color{kbvgrey}$=$} & \kbvresult{ಶ್ವೇತಾಂಬರ}{\color{kbvgrey} {[}ವಿ.ಪೂ.ಕ{]}} \\
\end{tabular}}
\end{adjustbox}\end{center}
\vspace{4pt plus 3pt minus 1pt}
\end{kbvderiv}

\indent {[}ಸೂಚನೆ : ಈ ಸಮಾಸಗಳಲ್ಲಿ ಪೂರ್ವಪದಗಳು ವಿಶೇಷಣಗಳಾಗಿಯೂ ಉತ್ತರಪದಗಳು ವಿಶೇಷ್ಯಗಳಾಗಿಯೂ ಇರುವುದರಿಂದ, ಇವುಗಳಿಗೆ ವಿಶೇಷಣ ಪೂರ್ವಪದ ಕರ್ಮಧಾರಯ ಸಮಾಸ {[}ವಿ.ಪೂ.ಕ{]} ವೆಂದು ಹೇಳುವುದು ವಾಡಿಕೆ.{]}\par

\kbvsechead{(೩)\hspace{0.45em}ದ್ವಿಗು ಸಮಾಸ :}

\indent - ಮೂರು {[}ಪೂ.ಪ{]} + ಲೋಕ {[}ಉ.ಪ{]} = ಮೂಲೋಕ ಇಲ್ಲಿ “ಮೂರು” ಎಂಬ ಸಂಖ್ಯಾವಾಚಕವು “ಲೋಕ” ಎಂಬ ನಾಮಪದ ದೊಂದಿಗೆ ಸೇರಿ ಸಮಸ್ತಪದವಾಗಿರುವುದು - ಇದು ದ್ವಿಗು ಸಮಾಸ ಎನಿಸುವುದು.\par
